\documentclass[11pt]{article}
\usepackage[margin=1in]{geometry}
\usepackage{amsmath, amssymb, amsthm}
\usepackage{hyperref}

\newcommand{\blank}{\underline{\hspace{1.5cm}}}

\title{MAT 3150 --- Homework 3\\
\large Problem 2 Walkthrough: Computing a Limit from the Definition}
\author{}
\date{}

\begin{document}
\maketitle

\textbf{Goal.} Use the limit definition to compute
\[
  \lim_{n\to\infty}\frac{n+2}{3n+5}.
\]

% ============================================================
\section*{The definition, in words}
$a_n \to L$ means:
\begin{quote}
For every $\varepsilon>0$ there is $N\in\mathbb{N}$ such that
$|a_n-L|<\varepsilon$ for all $n\ge N$.
\end{quote}
Think of it as a game. Someone hands you a tolerance $\varepsilon$
(maybe tiny, like $0.001$). You must answer with a point $N$ after which
\emph{every} term is within $\varepsilon$ of $L$. A proof is a recipe for
$N$ that works no matter which $\varepsilon$ you are handed.

\textbf{Tool you'll need.} Archimedean property: for any real number
$r$, there is a natural number $N$ with $N>r$.

\textbf{The plan (every problem of this type):}
\begin{enumerate}
  \item Guess $L$.
  \item Scratch work: simplify $|a_n-L|$, bound it by something simple,
        and solve ``simple thing $<\varepsilon$'' for $n$.
  \item Write the proof forward: ``Let $\varepsilon>0$. Choose $N$\ldots''
\end{enumerate}

% ============================================================
\section*{Worked example: $\displaystyle\lim_{n\to\infty}\frac{2n+1}{n+3}$}

\paragraph{1. Guess.} For large $n$, the $+1$ and $+3$ barely matter, so
$\dfrac{2n+1}{n+3}\approx\dfrac{2n}{n}=2$. Guess $L=2$.

\paragraph{2. Scratch work.} Common denominator:
\[
  \frac{2n+1}{n+3}-2
  = \frac{(2n+1)-2(n+3)}{n+3}
  = \frac{2n+1-2n-6}{n+3}
  = \frac{-5}{n+3}.
\]
The $n$'s in the numerator cancelled --- a good sign the guess is right.
So
\[
  \left|\frac{2n+1}{n+3}-2\right| = \frac{5}{n+3}.
\]
Simplify by making the denominator \emph{smaller}, which makes the
fraction \emph{bigger} (an overestimate is fine):
\[
  \frac{5}{n+3} < \frac{5}{n}.
\]
Now solve $\dfrac{5}{n}<\varepsilon$ for $n$: \quad
$n>\dfrac{5}{\varepsilon}$. That tells us how big $N$ must be.

\paragraph{3. The proof (what you hand in).}
\begin{proof}
We claim the limit is $2$. Let $\varepsilon>0$. By the Archimedean
property, choose $N\in\mathbb{N}$ with $N>\dfrac{5}{\varepsilon}$. Then
for all $n\ge N$,
\[
  \left|\frac{2n+1}{n+3}-2\right|
  = \frac{5}{n+3}
  < \frac{5}{n}
  \le \frac{5}{N}
  < \varepsilon .
\]
Hence $\displaystyle\lim_{n\to\infty}\frac{2n+1}{n+3}=2$.
\end{proof}

\noindent\textbf{Why each step in the chain holds:}
\begin{itemize}
  \item $=\,\dfrac{5}{n+3}$: the scratch-work algebra.
  \item $<\,\dfrac{5}{n}$: because $n+3>n$.
  \item $\le\,\dfrac{5}{N}$: because $n\ge N$.
  \item $<\,\varepsilon$: because $N>\dfrac{5}{\varepsilon}$, i.e.\
        $\dfrac{5}{N}<\varepsilon$.
\end{itemize}

% ============================================================
\newpage
\section*{Your turn: $\displaystyle\lim_{n\to\infty}\frac{n+2}{3n+5}$}

\paragraph{1. Guess.} For large $n$,
$\dfrac{n+2}{3n+5}\approx\dfrac{n}{3n}=\blank$. So $L=\blank$.

\paragraph{2. Scratch work.} Put it over a common denominator:
\[
  \frac{n+2}{3n+5}-L = \frac{\blank}{\blank}.
\]
\emph{Tip:} if $L$ is a fraction, multiply out carefully. If your guess
is right, the $n$'s in the numerator should cancel, leaving a constant.

\vspace{2cm}

Take absolute values:
\[
  \left|\frac{n+2}{3n+5}-L\right| = \blank .
\]

Now make the denominator smaller to get something simpler. (What is
$3n+5$ bigger than? You may need to handle an extra constant factor in
the denominator too.)
\[
  \blank \;<\; \blank .
\]

Solve ``simple bound $<\varepsilon$'' for $n$:
\[
  n > \blank .
\]

\vspace{1cm}

\paragraph{3. The proof.} Fill in the template:
\begin{quote}
We claim the limit is \blank. Let $\varepsilon>0$. By the Archimedean
property, choose $N\in\mathbb{N}$ with $N>\blank$. Then for all
$n\ge N$,
\[
  \left|\frac{n+2}{3n+5}-\blank\right|
  = \blank
  < \blank
  \le \blank
  < \varepsilon .
\]
Hence $\displaystyle\lim_{n\to\infty}\frac{n+2}{3n+5}=\blank$.
\end{quote}

\paragraph{Check yourself.} Can you justify each ``$=$'', ``$<$'',
``$\le$'' in your chain the same way the worked example does?

\end{document}
